% SEGMENT OLP-0009-S01
% Part: sets-functions-relations
% Chapter: sets
% Section: pairs-and-products

\documentclass[../../../include/open-logic-section]{subfiles}

\begin{document}

\olfileid{sfr}{set}{pai}
\olsection{ஜோடிகள், வரிசைத் தொகுதிகள், கார்டீசியன் பெருக்கங்கள்}

\begin{explain}
கணங்களின் உறுப்புகளுக்கு வரிசை இல்லை என்பது உறுப்புசார் சமத்துவத்திலிருந்து
பின்வருகிறது. எனவே வரிசையைக் குறிப்பிட விரும்பினால்,
$\tuple{x, y}$ என்னும் \emph{வரிசைப்படுத்தப்பட்ட ஜோடிகளைப்} பயன்படுத்துகிறோம்.
$\{x, y\}$ என்னும் வரிசைப்படுத்தப்படாத ஜோடியில் வரிசை பொருட்டல்ல:
$\{x, y\} = \{y, x\}$. வரிசைப்படுத்தப்பட்ட ஜோடியில் அது பொருட்டாகும்:
$x \neq y$ எனில் $\tuple{x, y} \neq \tuple{y, x}$.

கணக்கோட்பாட்டில் வரிசைப்படுத்தப்பட்ட ஜோடிகளை எவ்வாறு கருதுவது?
அவற்றின் முதல் உறுப்புகள் ஒன்றாகவும் இரண்டாம் உறுப்புகள் ஒன்றாகவும்
இருக்கும்போது, அப்போதும் அப்போது மட்டுமே அந்த ஜோடிகள் ஒன்றே என்ற
கருத்தைக் காக்க வேண்டும். அதாவது:
\[
  \tuple{a, b}= \tuple{c, d}\text{ அப்போதும் அப்போது மட்டுமே }a = c \text{ மற்றும் }b=d.
\]
வீனர்--குரடோவ்ஸ்கி வரையறையைப் பயன்படுத்திக் கணக்கோட்பாட்டில்
வரிசைப்படுத்தப்பட்ட ஜோடிகளை வரையறுக்கலாம்.
\end{explain}

% SEGMENT OLP-0009-S02
\begin{defn}[வரிசைப்படுத்தப்பட்ட ஜோடி]\ollabel{wienerkuratowski}
	$\tuple{a, b} = \{\{a\}, \{a, b\}\}$.
\end{defn}

% SEGMENT OLP-0009-S03
\begin{prob}
\olref[sfr][set][pai]{wienerkuratowski} ஐப் பயன்படுத்தி,
$a = c$, $b=d$ ஆகிய இரண்டும் நிறைவேறும்போது, அப்போதும் அப்போது மட்டுமே
$\tuple{a, b}= \tuple{c, d}$ என நிறுவுக.
\end{prob}

% SEGMENT OLP-0009-S04
\begin{explain}
வரிசைப்படுத்தப்பட்ட ஜோடிக்கு ஒரு வரையறையை நிர்ணயித்த பிறகு,
அதனைக் கொண்டு மேலும் கணங்களை வரையறுக்கலாம்.
எடுத்துக்காட்டாக, சில சமயங்களில் இரண்டு பொருள்களுக்கு மேற்பட்ட
வரிசைப்படுத்தப்பட்ட தொடர்முறைகளும் தேவைப்படும்:
$\tuple{x, y, z}$ என்னும் \emph{மூன்று-தொகுதிகள்},
$\tuple{x, y, z, u}$ என்னும் \emph{நான்கு-தொகுதிகள்} போன்றவை.
மூன்று-தொகுதிகளைச் சிறப்புவகை வரிசைப்படுத்தப்பட்ட ஜோடிகளாகக் கருதலாம்;
அவற்றின் முதல் உறுப்பே ஒரு வரிசைப்படுத்தப்பட்ட ஜோடியாக இருக்கும்:
$\tuple{x, y, z}$ என்பது $\tuple{\tuple{x, y},z}$ ஆகும்.
நான்கு-தொகுதிகளுக்கும் இதுவே பொருந்தும்:
$\tuple{x,y,z,u}$ என்பது $\tuple{\tuple{\tuple{x,y},z},u}$ ஆகும்;
இப்படியே தொடரலாம். பொதுவாக, $\tuple{x_1, \dots, x_n}$ என்னும்
\emph{வரிசைப்படுத்தப்பட்ட $n$-தொகுதிகள்} பற்றிப் பேசுகிறோம்.

வரிசைப்படுத்தப்பட்ட ஜோடிகளையோ பிற வரிசைப்படுத்தப்பட்ட
$n$-தொகுதிகளையோ கொண்ட சில கணங்கள் பயனுள்ளதாக இருக்கும்.
\end{explain}

% SEGMENT OLP-0009-S05
\begin{defn}[கார்டீசியன் பெருக்கம்]
$A$, $B$ ஆகிய கணங்கள் தரப்பட்டால், அவற்றின் \emph{கார்டீசியன் பெருக்கமான}
$A \times B$ பின்வருமாறு வரையறுக்கப்படுகிறது:
\[
  A \times B = \Setabs{\tuple{x, y}}{x \in A \text{ மற்றும் } y \in B}.
\]
\end{defn}

% SEGMENT OLP-0009-S06
\begin{ex}
$A = \{0, 1\}$, $B = \{1, a, b\}$ எனில், அவற்றின் பெருக்கம்
\[
A \times B = \{ \tuple{0, 1}, \tuple{0, a}, \tuple{0, b},
    \tuple{1, 1}, \tuple{1, a}, \tuple{1, b} \}.
\]
\end{ex}

% SEGMENT OLP-0009-S07
\begin{ex}
$A$ ஒரு கணம் எனில், $A$ ஐ அதனுடனேயே பெருக்கும் $A \times A$
என்பதை $A^2$ என்றும் எழுதுகிறோம். இது $x, y \in A$ என்ற நிபந்தனையை
நிறைவேற்றும் $\tuple{x, y}$ என்னும் \emph{அனைத்து} ஜோடிகளின் கணமாகும்.
$\tuple{x, y, z}$ என்னும் அனைத்து மூன்று-தொகுதிகளின் கணம் $A^3$ ஆகும்;
இப்படியே தொடரலாம். ஒரு மறுநிகழ்வு வரையறையைத் தரலாம்:
\begin{align*}
  A^1 & = A\\
  A^{k+1} & = A^k \times A
\end{align*}
\end{ex}

% SEGMENT OLP-0009-S08
\begin{prob}
$\{1, 2, 3\}^3$ இன் அனைத்து !!{element}s[acc] பட்டியலிடுக.
\end{prob}

% SEGMENT OLP-0009-S09
\begin{prop}\ollabel{cardnmprod}
$A$ இல் $n$ !!{element}s[um] $B$ இல் $m$ !!{element}s[um] இருந்தால்,
$A \times B$ இல் $n\cdot m$ உறுப்புகள் இருக்கும்.
\end{prop}

% SEGMENT OLP-0009-S10
\begin{proof}
$A$ இல் உள்ள ஒவ்வோர் !!{element} $x$ க்கும்,
$\tuple{x, y} \in A \times B$ என்ற வடிவில் $m$ !!{element}s உள்ளன.
$B_x = \Setabs{\tuple{x, y}}{y \in B}$ என்க.
$x_1 \neq x_2$ எனும்போதெல்லாம் $\tuple{x_1, y} \neq \tuple{x_2, y}$
என்பதால், $B_{x_1} \cap B_{x_2} = \emptyset$.
ஆனால் $A = \{x_1, \dots, x_n\}$ எனில்,
$A \times B = B_{x_1} \cup \dots \cup B_{x_n}$ ஆகும்;
ஆகவே அதில் $n\cdot m$ !!{element}s உள்ளன.

% SEGMENT OLP-0009-S11
இதனைக் காட்சிப்படுத்த, $A \times B$ இன் !!{element}s[acc]
ஒரு கட்ட அமைப்பில் வரிசைப்படுத்துக:
\[
\begin{array}{rcccc}
  B_{x_1} = & \{\tuple{x_1, y_1} & \tuple{x_1, y_2} & \dots & \tuple{x_1, y_m}\}\\
  B_{x_2} = & \{\tuple{x_2, y_1} & \tuple{x_2, y_2} & \dots & \tuple{x_2, y_m}\}\\
  \vdots & & \vdots\\
  B_{x_n} = & \{\tuple{x_n, y_1} & \tuple{x_n, y_2} & \dots & \tuple{x_n, y_m}\}
\end{array}
\]
$x_i$ அனைத்தும் வெவ்வேறானவை; $y_j$ அனைத்தும் வெவ்வேறானவை.
எனவே இந்தக் கட்ட அமைப்பில் எந்த இரு ஜோடிகளும் ஒன்றாக இல்லை;
அவற்றின் எண்ணிக்கை $n\cdot m$ ஆகும்.
\end{proof}

% SEGMENT OLP-0009-S12
\begin{prob}
$k$ மீதான தொகுத்தறிதலைப் பயன்படுத்தி, அனைத்து $k \ge 1$ க்கும்,
$A$ இல் $n$ !!{element}s இருந்தால் $A^k$ இல் $n^k$ !!{element}s
இருக்கும் எனக் காட்டுக.
\end{prob}

% SEGMENT OLP-0009-S13
\begin{ex}
$A$ ஒரு கணம் எனில், $A$ இன் !!{element}s கொண்டு அமைந்த எந்தத் தொடர்முறையும்
$A$ இன் மீதான ஒரு \emph{சொல்} எனப்படும்.
ஒரு தொடர்முறையை $A$ இன் !!{element}s கொண்டு அமைந்த $n$-தொகுதியாகக் கருதலாம்.
எடுத்துக்காட்டாக, $A = \{a, b, c\}$ எனில்,
``$bac$'' என்னும் தொடர்முறையை $\tuple{b, a, c}$ என்னும் மூன்று-தொகுதியாகக் கருதலாம்.
சொற்கள், அதாவது குறியீடுகளின் தொடர்முறைகள், கணினி அறிவியலில் மிக முக்கியமானவை.
வழக்கப்படி, $A$ இன் !!{element}s[acc] நீளம் $1$ உடைய தொடர்முறைகளாகவும்,
$\emptyset$ ஐ நீளம் $0$ உடைய தொடர்முறையாகவும் எண்ணுகிறோம்.
அப்போது $A$ இன் மீதான \emph{அனைத்துச்} சொற்களின் கணம்
\[
A^* = \{\emptyset\} \cup A \cup A^2 \cup A^3 \cup \dots
\]
ஆகும்.
\end{ex}

\end{document}
